\section{Additive valuations, discrete valuations, and DVRs}

\begin{definition}[Valuation ring and valuation ideal]
Let \((K,v)\) be a valued field. Define
\[
    R_v := \{a\in K \mid v(a)\geq 0\}
\]
and
\[
    P_v := \{a\in K \mid v(a)>0\}.
\]
Then \(R_v\) is called the \emph{valuation ring} of \(v\), and \(P_v\) is called the
\emph{valuation ideal} of \(v\).
\end{definition}

\begin{proposition}
Let \((K,v)\) be a valued field. Then:
\begin{enumerate}
    \item \(R_v\) is a subring of \(K\).
    \item \(P_v\) is an ideal of \(R_v\).
    \item The group of units of \(R_v\) is
    \[
        R_v^\times=\{a\in K^\times \mid v(a)=0\}.
    \]
    \item One has
    \[
        K=R_v\cup R_v^{-1},
    \]
    where
    \[
        R_v^{-1}:=\{a^{-1}\mid a\in R_v,\ a\neq 0\}.
    \]
    In particular, \(K\) is the field of fractions of \(R_v\).
    \item \(R_v\) is a local ring and
    \[
        J(R_v)=P_v.
    \]
    Equivalently, \(P_v\) is precisely the set of non-units of \(R_v\).
    \item The quotient field
    \[
        \kappa(v):=R_v/P_v
    \]
    is called the \emph{residue field} of \(v\).
\end{enumerate}
\end{proposition}

\begin{proof}
First, \(0,1\in R_v\). If \(a,b\in R_v\), then
\[
    v(ab)=v(a)+v(b)\geq 0,
\]
so \(ab\in R_v\). Also,
\[
    v(a+b)\geq \min\{v(a),v(b)\}\geq 0,
\]
so \(a+b\in R_v\). Finally, since \(v(-a)=v(a)\), we have \(-a\in R_v\).
Therefore \(R_v\) is a subring of \(K\).

Now let \(a,b\in P_v\). Then
\[
    v(a+b)\geq \min\{v(a),v(b)\}>0,
\]
so \(a+b\in P_v\). If \(r\in R_v\) and \(a\in P_v\), then
\[
    v(ra)=v(r)+v(a)>0.
\]
Thus \(ra\in P_v\), and hence \(P_v\) is an ideal of \(R_v\).

Next, let \(a\in R_v\). If \(a\) is a unit of \(R_v\), then \(a^{-1}\in R_v\), so
\[
    0=v(1)=v(aa^{-1})=v(a)+v(a^{-1}).
\]
Since \(v(a)\geq 0\) and \(v(a^{-1})\geq 0\), we get \(v(a)=0\). Conversely, if
\(v(a)=0\), then
\[
    v(a^{-1})=-v(a)=0,
\]
so \(a^{-1}\in R_v\). Therefore
\[
    R_v^\times=\{a\in K^\times\mid v(a)=0\}.
\]

For every \(a\in K^\times\), either \(v(a)\geq 0\), in which case \(a\in R_v\), or
\(v(a)<0\), in which case
\[
    v(a^{-1})=-v(a)>0,
\]
so \(a^{-1}\in R_v\). Hence
\[
    K=R_v\cup R_v^{-1}.
\]
In particular, every element of \(K\) is a fraction of elements of \(R_v\), so
\[
    K=\operatorname{Frac}(R_v).
\]

Finally, the non-units of \(R_v\) are exactly the elements \(a\in R_v\) satisfying
\(v(a)>0\), namely the elements of \(P_v\). Hence \(R_v\) has a unique maximal ideal
\(P_v\), so \(R_v\) is local and
\[
    J(R_v)=P_v.
\]
Consequently,
\[
    \kappa(v):=R_v/P_v
\]
is a field.
\end{proof}

\begin{definition}[Equivalent valuations]
Let \(v\) be a valuation on \(K\), and let \(\lambda\in\mathbb{R}_{>0}\). Define
\[
    v_\lambda(a):=\lambda v(a),\qquad a\in K^\times.
\]
Then \(v_\lambda\) is again a valuation on \(K\). We say that \(v\) is
\emph{equivalent} to \(v_\lambda\), and write
\[
    v\sim v_\lambda.
\]
\end{definition}

\begin{proposition}
Equivalent valuations have the same valuation ring, valuation ideal, group of units,
and residue field.
\end{proposition}

\begin{proof}
Let \(v_\lambda=\lambda v\), where \(\lambda>0\). Then for \(a\in K^\times\),
\[
    v_\lambda(a)\geq 0
    \iff
    \lambda v(a)\geq 0
    \iff
    v(a)\geq 0.
\]
Hence
\[
    R_{v_\lambda}=R_v.
\]
Similarly,
\[
    v_\lambda(a)>0
    \iff
    v(a)>0,
\]
so
\[
    P_{v_\lambda}=P_v.
\]
Also,
\[
    v_\lambda(a)=0
    \iff
    v(a)=0,
\]
and therefore
\[
    R_{v_\lambda}^\times=R_v^\times.
\]
Since the valuation rings and valuation ideals are the same, the residue fields are also
the same:
\[
    R_{v_\lambda}/P_{v_\lambda}=R_v/P_v.
\]
\end{proof}

\begin{definition}[Discrete valuation]
A valuation \(v\) of \(K\) is called \emph{discrete} if its value group is isomorphic to
\(\mathbb{Z}\), that is,
\[
    v(K^\times)\cong \mathbb{Z}.
\]
If moreover
\[
    v(K^\times)=\mathbb{Z},
\]
then \(v\) is called a \emph{normalized discrete valuation}.
\end{definition}

\begin{remark}
If \(v\) is a discrete valuation, then after replacing \(v\) by an equivalent valuation
\(\lambda v\) for a suitable \(\lambda>0\), one may assume that
\[
    v(K^\times)=\mathbb{Z}.
\]
Thus every discrete valuation is equivalent to a normalized discrete valuation.
\end{remark}

\begin{definition} 
    An element \(\pi\in R\) satisfying
    \[
        v(\pi)=1
    \]
    is called a uniformizer.
\end{definition}

\begin{theorem}
    Let \((K,v)\) be a field with a normalized discrete valuation. Let
    \[
        R=\{x\in K\mid v(x)\geq 0\}
    \]
    be its valuation ring, and let
    \[
        P=\{x\in R\mid v(x)>0\}
    \]
    be its maximal ideal. Suppose \(\pi\in R\) satisfies \(v(\pi)=1\). Then:
    \begin{enumerate}
        \item For every \(\alpha\in K^\times\), one can write
              \[
                  \alpha=\pi^{v(\alpha)}u
              \]
              for some unit \(u\in R^\times\).

        \item Every nonzero ideal \(I\) of \(R\) is of the form
              \[
                  I=(\pi^i)
              \]
              for some integer \(i\geq 0\). In particular,
              \[
                  P=(\pi),
                  \qquad
                  I=P^i,
              \]
              In particular, \(R\) is a principal ideal domain.

        \item The ring \(R\) is integrally closed.
    \end{enumerate}
\end{theorem}

\begin{proof}
    \textbf{(1).}
    Let \(\alpha\in K^\times\), and set
    \[
        i=v(\alpha).
    \]
    Since \(v(\pi)=1\), we have
    \[
        v(\pi^{-i}\alpha)
        =
        v(\alpha)-i v(\pi)
        =
        i-i
        =
        0.
    \]
    Therefore
    \[
        u:=\pi^{-i}\alpha\in R^\times.
    \]
    Hence
    \[
        \alpha=\pi^i u=\pi^{v(\alpha)}u,
    \]
    as desired.

    \medskip

    \textbf{(2).}
    Let \(I\neq 0\) be an ideal of \(R\). Since \(I\subseteq R\), for every nonzero
    \(a\in I\) we have
    \[
        v(a)\geq 0.
    \]
    Thus the set
    \[
        \{v(a)\mid 0\neq a\in I\}\subseteq \mathbb{Z}_{\geq 0}
    \]
    is nonempty, and therefore has a minimal element. Let
    \[
        i=\min\{v(a)\mid 0\neq a\in I\}.
    \]
    Choose \(a_0\in I\) such that
    \[
        v(a_0)=i.
    \]
    By part \((1)\), we may write
    \[
        a_0=\pi^i u
    \]
    for some \(u\in R^\times\). Hence
    \[
        \pi^i=a_0u^{-1}\in I.
    \]
    Therefore
    \[
        (\pi^i)\subseteq I.
    \]

    Conversely, let \(a\in I\) be nonzero. Write
    \[
        v(a)=j.
    \]
    By the minimality of \(i\), we have \(j\geq i\). Again by part \((1)\), there exists
    \(u\in R^\times\) such that
    \[
        a=\pi^j u.
    \]
    Then
    \[
        a
        =
        \pi^i\bigl(\pi^{j-i}u\bigr).
    \]
    Since \(j-i\geq 0\), we have \(\pi^{j-i}u\in R\). Therefore
    \[
        a\in (\pi^i).
    \]
    Hence
    \[
        I\subseteq (\pi^i).
    \]
    Combining the two inclusions, we obtain
    \[
        I=(\pi^i).
    \]

    In particular, applying this to the maximal ideal \(P\), we get
    \[
        P=(\pi).
    \]
    Therefore every nonzero ideal of \(R\) is of the form
    \[
        I=(\pi^i)=P^i.
    \]
    Since the zero ideal is also principal, \(R\) is a principal ideal domain.

    \medskip

    \textbf{(3).}
    By part \((2)\), the ring \(R\) is a principal ideal domain. Hence \(R\) is a unique
    factorization domain. Since every unique factorization domain is integrally closed,
    it follows that \(R\) is integrally closed.
\end{proof}

\begin{corollary}
    Let \(v\) and \(v'\) be two discrete valuations on a field \(K\). Then \(v\)
    and \(v'\) are equivalent if and only if they have the same valuation ideal, that is,
    \[
        P_v=P_{v'},
    \]
    where
    \[
        P_v=\{x\in K\mid v(x)>0\},
        \qquad
        P_{v'}=\{x\in K\mid v'(x)>0\}.
    \]
\end{corollary}

\begin{proof}
    If \(v\) and \(v'\) are equivalent, then there exists \(\lambda\in \mathbb{R}_{>0}\)
    such that
    \[
        v'=\lambda v.
    \]
    Hence, for every \(x\in K\),
    \[
        v'(x)>0 \Longleftrightarrow \lambda v(x)>0
        \Longleftrightarrow v(x)>0.
    \]
    Therefore
    \[
        P_v=P_{v'}.
    \]

    Conversely, suppose that
    \[
        P_v=P_{v'}=:P.
    \]
    We first show that the valuation ring can be recovered from the valuation ideal. Namely,
    \[
        R_v=K\setminus P^{-1},
    \]
    where
    \[
        P^{-1}:=\{a^{-1}\mid a\in P,\ a\neq 0\}.
    \]
    Indeed, if \(x\in R_v\), then \(v(x)\geq 0\). Hence \(x\) cannot be of the form
    \(a^{-1}\) with \(a\in P\), because then
    \[
        v(x)=v(a^{-1})=-v(a)<0.
    \]
    Thus \(x\notin P^{-1}\). Conversely, if \(x\notin R_v\), then \(x\neq 0\) and
    \(v(x)<0\). Hence
    \[
        v(x^{-1})=-v(x)>0,
    \]
    so \(x^{-1}\in P\), and therefore
    \[
        x=(x^{-1})^{-1}\in P^{-1}.
    \]
    This proves
    \[
        R_v=K\setminus P^{-1}.
    \]
    Since \(P_v=P_{v'}\), we get
    \[
        R_v=R_{v'}=:R.
    \]

    Now replace \(v\) and \(v'\) by equivalent normalized discrete valuations. This does
    not change the valuation ring or the valuation ideal. Hence we may assume that
    \[
        v(K^\times)=\mathbb{Z},
        \qquad
        v'(K^\times)=\mathbb{Z}.
    \]

    Choose \(\pi\in K^\times\) such that
    \[
        v(\pi)=1.
    \]
    Then \(\pi\) is a uniformizer for \(v\), so by Theorem~9.0.11,
    \[
        P=(\pi)
    \]
    as an ideal of \(R\).

    Let \(\pi'\) be a uniformizer for \(v'\). Then similarly
    \[
        P=(\pi').
    \]
    Since
    \[
        (\pi)=(\pi'),
    \]
    there exist \(r,s\in R\) such that
    \[
        \pi=r\pi',
        \qquad
        \pi'=s\pi.
    \]
    Applying \(v'\), we obtain
    \[
        v'(\pi)=v'(r)+v'(\pi')\geq 1,
    \]
    and
    \[
        1=v'(\pi')=v'(s)+v'(\pi)\geq v'(\pi).
    \]
    Therefore
    \[
        v'(\pi)=1.
    \]

    Now let \(a\in K^\times\). By Theorem~9.0.11, there exists \(u\in R^\times\) such
    that
    \[
        a=\pi^{v(a)}u.
    \]
    Since \(R=R_{v'}\), the units of \(R\) are exactly the elements of valuation \(0\) with
    respect to \(v'\). Hence
    \[
        v'(u)=0.
    \]
    Therefore
    \[
        v'(a)=v\left(a\right)v'(\pi)+v'(u)=v(a).
    \]
    Thus \(v'=v\) after normalization.

    Hence the original discrete valuations differ by multiplication by a positive real
    constant. Therefore \(v\) and \(v'\) are equivalent.
\end{proof}

\begin{theorem}[Discrete valuation rings and local PIDs]
    Let \(R\) be an integral domain which is not a field, and let
    \(K=\operatorname{Frac}(R)\). Then \(R\) is a discrete valuation ring if and only if
    \(R\) is a local principal ideal domain.
\end{theorem}

\begin{proof}
    Suppose first that \(R\) is a discrete valuation ring. Then there exists a normalized
    discrete valuation
    \[
        v:K^\times\longrightarrow \mathbb{Z}
    \]
    such that
    \[
        R=\{x\in K\mid v(x)\geq 0\}.
    \]
    By the structure theorem for valuation rings of normalized discrete valuations, \(R\)
    is a local principal ideal domain. More precisely, if \(\pi\in R\) satisfies
    \(v(\pi)=1\), then the maximal ideal of \(R\) is
    \[
        P=\{x\in R\mid v(x)>0\}=(\pi),
    \]
    and every nonzero ideal of \(R\) is of the form
    \[
        P^n=(\pi^n)
    \]
    for some integer \(n\geq 0\). Hence \(R\) is a local PID.

    Conversely, suppose that \(R\) is a local principal ideal domain. Since \(R\) is not
    a field, its unique maximal ideal is nonzero. Denote it by
    \[
        P=J(R).
    \]
    Since \(R\) is a PID, there exists a nonzero nonunit \(\pi\in R\) such that
    \[
        P=(\pi).
    \]

    We first claim that
    \[
        \bigcap_{n\geq 1}P^n=0.
    \]
    Suppose, for contradiction, that
    \[
        0\neq a\in \bigcap_{n\geq 1}P^n.
    \]
    Since \(P^n=(\pi^n)\), for each \(n\geq 1\) there exists \(a_n\in R\) such that
    \[
        a=\pi^n a_n.
    \]
    Then
    \[
        \pi^n a_n=\pi^{n+1}a_{n+1}.
    \]
    Since \(R\) is an integral domain and \(\pi\neq 0\), we may cancel \(\pi^n\), obtaining
    \[
        a_n=\pi a_{n+1}.
    \]
    Hence
    \[
        (a_n)\subseteq (a_{n+1}).
    \]
    This inclusion is strict. Indeed, if \((a_n)=(a_{n+1})\), then
    \(a_{n+1}=r a_n\) for some \(r\in R\). Hence
    \[
        a_{n+1}=r\pi a_{n+1}.
    \]
    Since \(a_{n+1}\neq 0\), we get \(1=r\pi\), which means that \(\pi\) is a unit, a
    contradiction. Thus we obtain a strictly ascending chain of ideals
    \[
        (a_1)\subsetneq (a_2)\subsetneq (a_3)\subsetneq\cdots.
    \]
    This contradicts the fact that \(R\) is Noetherian, because every PID is Noetherian.
    Therefore
    \[
        \bigcap_{n\geq 1}P^n=0.
    \]

    Now let \(0\neq a\in R\). Since \(P^0=R\), the set
    \[
        \{n\geq 0\mid a\in P^n\}
    \]
    is nonempty. By the claim above, \(a\) cannot belong to all powers \(P^n\). Hence
    there exists a unique integer \(n\geq 0\) such that
    \[
        a\in P^n
        \qquad\text{and}\qquad
        a\notin P^{n+1}.
    \]
    Since \(P^n=(\pi^n)\), we may write
    \[
        a=\pi^n u
    \]
    for some \(u\in R\). If \(u\in P\), then
    \[
        a=\pi^n u\in \pi^n P=P^{n+1},
    \]
    contradicting the choice of \(n\). Thus \(u\notin P\). Since \(R\) is local and \(P\)
    is the set of nonunits, it follows that
    \[
        u\in R^\times.
    \]
    Therefore every nonzero element \(a\in R\) can be written uniquely in the form
    \[
        a=\pi^n u,
        \qquad n\geq 0,\quad u\in R^\times.
    \]

    We now extend this to \(K^\times\). Every \(x\in K^\times\) can be written as
    \[
        x=\frac{a}{b},
        \qquad a,b\in R\setminus\{0\}.
    \]
    Write
    \[
        a=\pi^m u,\qquad b=\pi^n w,
    \]
    where \(m,n\geq 0\) and \(u,w\in R^\times\). Then
    \[
        x=\pi^{m-n}uw^{-1}.
    \]
    Hence every \(x\in K^\times\) has a representation
    \[
        x=\pi^r u,
        \qquad r\in \mathbb{Z},\quad u\in R^\times.
    \]
    This representation is unique. Indeed, if
    \[
        \pi^r u=\pi^s w,
        \qquad u,w\in R^\times,
    \]
    and say \(r<s\), then
    \[
        u=\pi^{s-r}w\in P,
    \]
    contradicting that \(u\) is a unit. Hence \(r=s\).

    Define
    \[
        v:K^\times\longrightarrow \mathbb{Z}
    \]
    by
    \[
        v(\pi^r u)=r,
        \qquad r\in\mathbb{Z},\quad u\in R^\times.
    \]
    This is well-defined by the uniqueness just proved.

    We verify that \(v\) is a valuation. If
    \[
        x=\pi^r u,\qquad y=\pi^s w,
    \]
    with \(u,w\in R^\times\), then
    \[
        xy=\pi^{r+s}uw,
    \]
    and \(uw\in R^\times\). Therefore
    \[
        v(xy)=r+s=v(x)+v(y).
    \]
    Next, suppose \(x+y\neq 0\). Without loss of generality, assume \(r\leq s\). Then
    \[
        x+y=\pi^r\left(u+\pi^{s-r}w\right).
    \]
    Since
    \[
        u+\pi^{s-r}w\in R,
    \]
    its valuation is nonnegative whenever it is nonzero. Thus
    \[
        v(x+y)\geq r=\min\{r,s\}=\min\{v(x),v(y)\}.
    \]
    Hence \(v\) is a valuation.

    Moreover,
    \[
        v(\pi)=1,
    \]
    so
    \[
        v(K^\times)=\mathbb{Z}.
    \]
    Therefore \(v\) is a normalized discrete valuation.

    Finally, we show that \(R\) is exactly the valuation ring of \(v\). Let
    \[
        R_v:=\{x\in K\mid v(x)\geq 0\}.
    \]
    If \(x\in R\setminus\{0\}\), then by the decomposition above,
    \[
        x=\pi^n u
    \]
    with \(n\geq 0\), so \(v(x)=n\geq 0\). Hence \(R\subseteq R_v\).

    Conversely, if \(x\in R_v\setminus\{0\}\), write
    \[
        x=\pi^r u,
        \qquad r\in\mathbb{Z},\quad u\in R^\times.
    \]
    Since \(v(x)=r\geq 0\), we have
    \[
        x=\pi^r u\in R.
    \]
    Hence \(R_v\subseteq R\).

    Therefore
    \[
        R=R_v.
    \]
    Thus \(R\) is the valuation ring of a normalized discrete valuation, and hence \(R\)
    is a discrete valuation ring.
\end{proof}
